\section{Conclusions and Future Scope}

\subsection{Work Accomplished}
Progress is compared with the four approved objectives. Each objective is split
into five reviewable milestones. Table~\ref{tab:objective_progress} reports
completed engineering work. The percentages are not accuracy scores.

\begin{table}[H]
\centering
\caption{Progress against approved objectives}
\label{tab:objective_progress}
\small
\begin{tabularx}{\textwidth}{p{0.8in}p{0.85in}X}
\toprule
\textbf{Objective} & \textbf{Coverage} & \textbf{Evidence and remaining check}\\
\midrule
1 & 4 of 5 & Vision, wearable, fusion, API, dashboard, and relay paths are present. Controlled end-to-end sync report remains.\\
2 & 4 of 5 & ESP32-S3 firmware, sensors, provisioning, and repeatable builds are present. Reference-device field test remains.\\
3 & 3 of 5 & Feature engine, tactical rules, and training path are present. Reviewed outcome data and model evaluation remain.\\
4 & 2 of 5 & Public-data training and software tests are present. Locked campus tracking, sensor, sync, and coach tests remain.\\
\bottomrule
\end{tabularx}
\end{table}

The milestone coverage is 80, 80, 60, and 40 percent. Current detector logs
show that the training pipeline runs. The 30-epoch record ends with validation
precision 0.8435, recall 0.7318, AP$_{50}$ 0.7476, and AP$_{50:95}$ 0.4589 on
its recorded split. These values are not a campus deployment result.

On August 20, 2026, all 53 project tests, frontend lint, and frontend production
build passed. The supported \texttt{wearable\_stream} firmware compiled for the
ESP32-S3 at 31 percent flash and 8 percent dynamic memory. Hardware reference
tests are still reported separately from compilation and software tests.

\subsection{Conclusions}
The mid-semester prototype gives a working path from camera and wearable data to
a live coaching view. The most important design choice is to keep identity,
time, and quality before calculating analytics. This makes wrong or weak values
easier to find and correct.

Main processing and evidence stay on the project PC. The required VPS relay
provides the public domain route for raw wearable telemetry and short
reconnection replay without storing a match archive. The tested firmware is the
main field build, and the model-training path is ready for reviewed campus
labels.

The literature also sets clear limits. Ball tracking and player identity are
difficult during occlusion. PPG needs movement tests, and workload models need
many sessions with consistent labels. The project therefore uses visible
quality states, manual review, component tests, and results linked to the
original session.

\subsection{Environmental, Economic, and Social Benefits}
Local video processing avoids sending full recordings to a cloud service. The
VPS keeps only a small memory cache, which also limits storage use. Replaceable
sensor modules and straps make repairs easier, while batteries and damaged
electronics still require proper disposal.

The design uses common components and open-source software. A controlled drill
can start with one or two wearables before increasing the count. This makes the
system suitable for a campus team, though time for labelling, fitting, and
validation remains part of the real cost.

The product can give coaches clearer evidence without requiring a high-budget
tracking system. Responsible use requires player consent, limited access, and
careful wording around health. A workload warning should support discussion
with the player and qualified staff, not label a player as injured.

\subsection{Future Work Plan}
Table~\ref{tab:future_plan} gives the remaining work and the evidence expected
from each stage.

\begingroup
\small
\begin{longtable}{@{}p{1.0in}p{2.15in}p{2.75in}@{}}
\caption{Future work plan}\label{tab:future_plan}\\
\toprule
\textbf{Period} & \textbf{Work} & \textbf{Acceptance evidence}\\
\midrule
\endfirsthead
\toprule
\textbf{Period} & \textbf{Work} & \textbf{Acceptance evidence}\\
\midrule
\endhead
Before panel & Add final screens and hardware photograph; rehearse the fixed demo. & Placeholders filled, ten-minute run, clean startup, relay reconnect.\\
August--September & Run mapping, tap-sync, sensor tests, and campus labelling. & Mapping error, 20 sync events, sensor error, reviewed labels, session split.\\
October & Train and test detector, ball, jersey, and identity updates. & Per-class AP, HOTA, IDF1, switch count, reviewed error clips.\\
October--November & Collect workload records and train simple fusion baselines. & Reviewed session labels, data gate, later-session test, Brier score.\\
Final phase & Run coach-use, failure, and privacy checks; complete the report. & Task notes, stated limits, signed evidence, repeatable build.\\
\bottomrule
\end{longtable}
\endgroup

Later improvements can include more cameras, better identity recovery after long
gaps, jersey recognition trained on local kits, and a higher-resolution ball
detector. Tactical work can compare passing options and short changes in team
shape. Each added result will keep its input quality, visible evidence, and
confidence.
